The improved Euler's method, also known as Heun's method, is a predictor-corrector approach that achieves better accuracy than the standard Euler method by using an average of slopes.
Using the same initial value problem as in ~\ref{subsubsec:euler-method}, the method proceeds in two sequential steps:

\[
y_{n+1}^* = y_n + f(t_n, y_n) \cdot \Delta t
\quad \rightarrow \quad
y_{n+1} = y_n + \frac{f(t_n, y_n) + f(t_{n+1}, y_{n+1}^*)}{2} \cdot \Delta t.
\]

The first step provides a predictor value using the standard Euler method.
This predicted value is then used to compute a more accurate slope at the next time point.
The second step takes the average of the slopes at the beginning and end of the interval to obtain the corrected value.

By using this averaging technique, the improved Euler method provides a better balance between computational efficiency and accuracy for many ordinary differential equations.